\section{Probelemstellung, Forschungsfrage und Vorgehen}
Das Erstellen von geeigneten Tests für Softwareprodukte stellt einen wichtigen Bestandteil der Qualitätssicherung dar. Mit ihnen wird überprüft, ob ein Programm
das vorab definierte Verhalten korrekt umsetzt und auch bei Sonderfällen zuverlässig bleibt und die erwarteten Ergebnisse liefert. Gerade bei kurzen algorithmischen 
Programmieraufgaben lässt sich häufig der Fall beobachten, dass eine falsche Implementierung für viele Standardfälle das korrekte Ergebnis liefert und nur 
bei einzelnen speziellen Eingaben/Tests fehlschlägt.

Da unterschiedliche fehlerhafte Lösungen zuverlässig erkannt werden müssen, stellt das Konzipieren geeigneter Testfälle im Bereich der Softwareentwicklung eine 
anspruchsvolle Aufgabe dar. Angesichts des zunehmenden Einsatzes von KI-Systemem im Feld der Softwareentwicklung entsteht die Frage, ob mit Hilfe der künstlichen Intelligenz 
nur Code erzeugt werden kann, oder ob die Modelle auch in der Lage sind geeignete Testfälle inklusive der erwarteten Ergebnisse, anhand von vorab definiertem Verhalten 
zu generieren.

Das Ziel dieser Projektarbeit besteht darin, ein KI-Assistenzsystem in Form einer Webanwengung zu entwickeln, welches Testfälle inklusive der erwarteten Ergebnisse 
für natürlichsprachliche Programmieraufgaben ausgibt. Auf Basis dieser Anwendung wird anschließend an 7 beispielhaften Aufgaben untersucht, wie valide die generierten 
Testfälle sind und wie gut sie fehlerhafte Implementierung einer Lösung erkennen. Wichtig hierbei ist, dass die KI keine Zugriff auf die Implementierungen selbst erhält und 
die Testfälle somit ausschließliche auf Basis des beschriebenden Sollverahltens innerhalb der Aufgabe abgeleitet werden. Daraus ergibt sich die folgende Forschungsfrage:

\begin{center}
    \textit{Inwiefern können Large Language Models aus natürlichsprachlichen Programmieraufgaben valide Black-Box-Testfälle generieren, die zur Erkennung fehlerhafter Implementierungen 
geeignet sind?}
\end{center}

Durch die Untersuchung soll herausgearbeitet werden, ob ein solches System als unterstützendes Werkzeug bei der Testfallerstellung eingesetzt werden kann und ob eine manuelle 
fachilche Prüfung der generierten Testfälle erforderlich bleibt.

Dafür werden zunächst die Grundlagen des Software-Testings und die Testfallgenerierung mithilfe von
KI erläutert. Dies schafft das fachliche Fundament für die spätere Bewertung der
generierten Testfälle. Anschließend werden die Konzeption und technische Umsetzung des
Assistenzsystems beschrieben, um zu herauszustellen, wie die theoretischen Anforderungen in einen
praktisch nutzbaren Ablauf übertragen wurden und aus welchen technischen Komponenten die Anwendung besteht. 
Darauf aufbauend werden die Evaluationsmethode sowie die Ergebnisse zur Validität und Fehlererkennung erläutert. 
Dieser Schritt ist wichtig, um die Forschungsfrage anhand nachvollziehbarer Kriterien beantworten zu können. 
Das abschließende Fazit ordnet die gewonnenen Erkenntnisse und ihre Grenzen ein und leitet daraus
mögliche Weiterentwicklungen ab.

\newpage